\subsection{推广到多模复形}

设 B 和 B' 为两个环，C 为 (B, A)-双模复形，C' 为 (A, B')-双模复形（X，第 43 页）；则 (C ⊗\_A C', D) （X，第 61 页）是 (B, B')-双模的复形，且典范同态

\[
\gamma : \mathrm{H} (\mathbf {C}) \otimes_ {\mathbf {A}} \mathrm{H} (\mathbf {C} ^ {\prime}) \rightarrow \mathrm{H} (\mathbf {C} \otimes_ {\mathbf {A}} \mathbf {C} ^ {\prime})
\]

与两个成员的(B, B')-双模结构相容。

若 B'' 是第三个环，且 C'' 是 (B', B'')-双模的复形，则典范同态（II，第 64 页，命题 8）

\[
\left(\mathrm{C} \otimes_ {\mathrm{A}} \mathrm{C} ^ {\prime}\right) \otimes_ {\mathrm{B} ^ {\prime}} \mathrm{C} ^ {\prime \prime} \rightarrow \mathrm{C} \otimes_ {\mathrm{A}} \left(\mathrm{C} ^ {\prime} \otimes_ {\mathrm{B} ^ {\prime \prime}} \mathrm{C} ^ {\prime \prime}\right)
\]

是(B, B'')-双模复形的一个同构。

更一般地，我们留给读者按照第 1 号（X，第 63 页）以及 II，第 65 至 72 页的模型，发展有限全序族多模复形的张量积理论（结合性与交换性同构， \ldots).

设 B 和 B′ 为两个环，M 为 (B, A)-双模，N 为 (A, B′)-双模：则 L(M) ⊗\_A L(N) 是 (B, B')-双模的复形，因此 Tor\^{}A (M, N) 具有自然的 (B, B')-双模分次结构；在 0 次项上，该结构与 M ⊗\_A N 的结构一致。

若 \(\lambda\in\mathbf{B},\lambda^{\prime}\in\mathbf{B}^{\prime}\) ，且若我们记 \(\lambda_{\mathbf{M}},\lambda_{\mathbf{N}}^{\prime},\lambda_{\mathbf{T}},\lambda_{\mathbf{T}}^{\prime}\) 为位似 \(x\mapsto\lambda x,y\mapsto y\lambda^{\prime},z\mapsto\lambda z,z\mapsto z\lambda^{\prime}\) ，它们分别作用于 \(\mathbf{M},\mathbf{N},\mathrm{Tor}^{\mathbf{A}}(\mathbf{M},\mathbf{N}),\mathrm{Tor}^{\mathbf{A}}(\mathbf{M},\mathbf{N})\) ，则

\[
\lambda_ {\mathrm{T}} = \operatorname{Tor} ^ {\mathbf {A}} \left(\lambda_ {\mathbf {M}}, 1 _ {\mathbf {N}}\right), \quad \lambda_ {\mathrm{T}} ^ {\prime} = \operatorname{Tor} ^ {\mathbf {A}} \left(1 _ {\mathbf {M}}, \lambda_ {\mathbf {N}} ^ {\prime}\right),
\]

这给出了Tor \(^{A}\) (M, N)双模结构的另一种描述。我们留给读者将第5号和第7号推广到多模复形的情形。

